\documentclass{beamer}

\usepackage{fullpage, parskip, graphicx, amsmath, amsfonts, amssymb, amsthm, xcolor, multicol}	
	\pagestyle{empty}
\usepackage{enumerate}

\title{Course Notes: Statistics 1401}
\subtitle{Final Exam Review}
\author{Robb Sinn}
\date{November 13, 2026}	

\usetheme[numbering=none]{focus}
\usecolortheme{seahorse}

\begin{document}
\frame{\titlepage}
\begin{frame}{Exam Review Questions}
Let's look at a few real-world examples of research studies and think about the results in terms of statistics.    
\end{frame}

\begin{frame}\frametitle{Example 1: The Trailhead Litter Contagion Study}
\textbf{Core Question:} Does the visible state of a trailhead—specifically the presence of a single piece of litter or social proof messaging—alter how hikers treat Leave No Trace principles miles into the backcountry?

\textbf{Methodology:} Select a popular hiking trail and randomly assign days to one of three trailhead conditions:
\begin{enumerate}
    \item Pristine trailhead with a sign reading ``94\% of hikers pack out all trash,"
    \item a single planted aluminum can next to the trailhead kiosk, or 
    \item a control condition with standard park signage.
\end{enumerate} Researchers stationed discreetly along the first mile and at a scenic overlook record littering behavior and whether hikers pick up planted trash.
\end{frame}

\begin{frame}{Hiking Question}
\textbf{What type of study is being conducted?}\\[9mm]
\pause
\textbf{Results:} If we found that $p = 0.0145$, what would our statistical conclusion be?\\[9mm]
\pause
\textbf{Reporting Out:} What would this mean in the real-world context of the question? 
\vspace{25mm}
\end{frame}

\begin{frame}\frametitle{Example 2: Paddling Pace Synchronization & Risk Tolerance}
\textbf{Core Question:} When novice paddlers join a group tour, do they subconsciously synchronize their stroke cadence and adopt higher risk tolerances (e.g., paddling closer to rapids or open water) based on the behavior of an undercover "pacemaker" paddler?

\textbf{Methodology:} Groups of recreational kayakers embark on a guided river or coastal paddle containing one confederate posing as a fellow tourist. 

\begin{enumerate}
    \item In Condition A, the confederate maintains an aggressive stroke rate and hugs the outer edge of the channel. 
    \item In Condition B, they paddle leisurely near the shore.
\end{enumerate} 
GPS trackers and GoPro mounts measure group cohesion, stroke cadence convergence, and distance from shoreline hazards.
\end{frame}

\begin{frame}{Kayaking Question}
If a single numeric measure of group cohesion (which includes cadence convergence and safety margins), \textbf{What type of study is being conducted?}\\[9mm]
\pause
\textbf{Results:} If we found that $p = 0.476$, what would our statistical conclusion be?\\[9mm]
\pause
\textbf{Reporting Out:} What would this mean in the real-world context of the question? 
\vspace{25mm}
\end{frame}

\begin{frame}\frametitle{Example 3. Campfire Storytelling \& Prosocial Bonding}
\textbf{Core Question:} Does the physical arrangement and lighting of an evening camp gathering (open fire vs. bright LED camp lanterns) shift group conversation toward self-disclosure, vulnerability, and cooperative camp chores the following morning?

\pause
\textbf{Methodology:} Randomly assign multi-day group camping trips to either traditional wood campfires or high-lumen white LED lanterns as their primary evening light source. Audio analysis (measuring conversational turn-taking and emotional valence) and morning observation of unprompted communal tasks (filtering water, washing dishes, packing shared gear) quantify social cohesion.

\pause
Researchers stationed discreetly along the first mile and at a scenic overlook record littering behavior and whether hikers pick up planted trash.
\end{frame}

\begin{frame}{Campfire Question}
\textbf{What type of study is being conducted?}\\[9mm]
\pause
\textbf{Results:} If we found that $p = 0.000393$, what would our statistical conclusion be?\\[9mm]
\pause
\textbf{Reporting Out:} What would this mean in the real-world context of the question? 
\vspace{25mm}
\end{frame}

\begin{frame}\frametitle{Example 4: The ``Drafting" \& Effort Contagion Effect}
\textbf{Core Question:} Does riding behind a cyclist displaying visible signs of fatigue versus effortless posture alter a trailing rider’s power output and perceived exertion at identical speeds?

\pause
\textbf{Methodology:} On a closed outdoor loop, recreational cyclists ride behind a confederate pacing at a fixed 18 mph (verified via power meter). In one condition, the confederate rides smoothly with relaxed shoulders; in the other, they rock their shoulders, breathe heavily, and glance down repeatedly. Trailing riders report Rating of Perceived Exertion (RPE) every lap while heart rate and cadence are logged.
\vspace{15mm}
\end{frame}

\begin{frame}{Cycling Question}
\textbf{What type of study is being conducted?}\\[9mm]
\pause
\textbf{Results:} If we found that $p = 0.05155$, what would our statistical conclusion be?\\[9mm]
\pause
\textbf{Reporting Out:} What would this mean in the real-world context of the question? 
\vspace{25mm}
\end{frame}

\begin{frame}\frametitle{Example 5. The ``Rabbit" Effect in Recreational Road Races}
\textbf{Core Question:} Can an unexpected surge by a non-elite runner in the middle of a pack trigger a localized "pace contagion" that causes nearby runners to abandon their pre-planned race pacing strategies?

\pause
\textbf{Methodology:} During a local 5K or 10K road race, embed pairs of confederates in mid-pack starting corrals wearing standard race bibs. At the halfway mark, confederates execute a synchronized 400-meter surge (dropping pace by 30 seconds/mile) before settling back. Chip-timing splits and wearable GPS uploads from surrounding runners reveal how far the pace wave ripples backward through the pack.
\vspace{15mm}
\end{frame}

\begin{frame}{Running Question}
\textbf{What type of study is being conducted?}\\[9mm]
\pause
\textbf{Results:} If we found that $p = 0.3317$, what would our statistical conclusion be?\\[9mm]
\pause
\textbf{Reporting Out:} What would this mean in the real-world context of the question? 
\vspace{25mm}
\end{frame}

\begin{frame}\frametitle{Example 6. Dining Hall Tray \& Portion Mimicry}
\textbf{Core Question:} Do college freshmen mirror the plate composition (vegetable-to-dessert ratio) and eating speed of the stranger sitting directly across from them in a campus dining hall?

\pause
\textbf{Methodology:} In a high-traffic university dining hall, a confederate sits at communal tables with either a heavily plant-based plate or a high-calorie comfort-food plate, eating at either a fast or slow bite cadence. Observers code the food selections and meal duration of students who randomly choose seats within two chairs of the confederate.
\vspace{15mm}
\end{frame}

\begin{frame}{Campus Dining Question}
If a single numeric variable measuring \textbf{dining conformity} were developed and recorded, \textbf{what type of study is being conducted?}\\[9mm]
\pause
\textbf{Results:} If we found that $p = 0.002$, what would our statistical conclusion be?\\[9mm]
\pause
\textbf{Reporting Out:} What would this mean in the real-world context of the question? 
\vspace{25mm}
\end{frame}

\begin{frame}\frametitle{Example 7. Library Focus vs. Distraction Cascades}
\textbf{Core Question:} In an open university library reading room, does the visible screen activity of neighboring students (academic work vs. streaming sports/social media) influence how quickly a newly arrived student switches away from their coursework?

\pause
\textbf{Methodology:} Confederates sit at central study tables with laptops angled slightly toward adjacent empty seats. In Condition A, their screens display dense PDFs and coding environments; in Condition B, they silently play YouTube videos or scroll feeds. Discreet observation logs the ``time-to-first-distraction" (picking up a phone or opening a non-academic browser tab) for students who sit beside them.

\vspace{15mm}
\end{frame}

\begin{frame}{Library Study Spaces Question}
The numeric variable is ``time-to-first distraction", so \textbf{what type of study is being conducted?}\\[9mm]
\pause
\textbf{Results:} If we found that $p = 0.009897$, what would our statistical conclusion be?\\[9mm]
\pause
\textbf{Reporting Out:} What would this mean in the real-world context of the question? 
\vspace{25mm}
\end{frame}

\begin{frame}{Example 8. The Standing Ovation Tipping Point}
\textbf{Core Question:} What critical mass of early standers is required to trigger a full-auditorium standing ovation, and does seating location (front row vs. aisle seats in the middle) change the contagion threshold?

\pause
\textbf{Methodology:} At campus performing arts or local theater productions, place small clusters of confederates (2, 4, or 6 people) in either the front center or mid-aisle rows. Regardless of performance quality, confederates stand and applaud enthusiastically within 3 seconds of the curtain call. Overhead balcony cameras track the wave propagation and total percentage of the audience that stands.
\vspace{23mm}
\end{frame}

\begin{frame}{Campus Recreation Questions}
\textbf{What type of study is being conducted?}\\[9mm]
\pause
\textbf{Results:} If we found that $p = 0.707$, what would our statistical conclusion be?\\[9mm]
\pause
\textbf{Reporting Out:} What would this mean in the real-world context of the question? 
\vspace{25mm}
\end{frame}

\begin{frame}{Example 9: The Trivia Night ``Confident Whisper" Conformity Test}
\textbf{Core Question:} In a noisy pub trivia setting, does overhearing a nearby table confidently state an incorrect answer cause neighboring teams to second-guess their own correct knowledge and switch their final submission?

\pause
\textbf{Methodology:} During a weekly bar trivia night, seat a team of confederates at a central table within earshot of surrounding tables. On three moderately difficult questions where most players know the right answer (e.g., ``What is the capital of Australia?"), the confederates loudly and confidently agree on a plausible wrong answer (``Sydney, definitely Sydney—I visited last year") right before answer sheets are collected. Researchers analyze answer-sheet cross-outs and final submissions by proximity to the confederate table.

\pause
\footnotesize{ \textbf{Note the classic Asch-style} informational conformity adapted to an informal, competitive social entertainment environment.}
\end{frame}

\begin{frame}{Trivia Night Questions}
\textbf{What type of study is being conducted?}\\[9mm]
\pause
\textbf{Results:} If we found that $p = 0.00107$, what would our statistical conclusion be?\\[9mm]
\pause
\textbf{Reporting Out:} What would this mean in the real-world context of the question? 
\vspace{25mm}
\end{frame}


\end{document}